%==============================================================================
\section{Hồi quy logistic (Logistic Regression)}
\label{sec:logistic-regression}
%==============================================================================

\begin{dinhnghia}[title={Giới thiệu hồi quy logistic}]
Hồi quy logistic là thuật toán phân loại học có giám sát dùng để dự đoán \textbf{xác suất} của biến mục tiêu.
Bản chất biến phụ thuộc (target) là \textbf{dichotomous} --- chỉ có hai lớp có thể.
\end{dinhnghia}

\begin{ynghia}[title={Biến nhị phân}]
Nói đơn giản, biến phụ thuộc có dạng nhị phân: dữ liệu mã hóa \textbf{1} (thành công/có) hoặc \textbf{0} (thất bại/không).
\end{ynghia}

\begin{dinhnghia}[title={Mô hình xác suất}]
Về mặt toán học, mô hình hồi quy logistic dự đoán $P(Y=1)$ như một hàm của $X$.
Đây là một trong những thuật toán ML đơn giản nhất, áp dụng cho nhiều bài toán phân loại: phát hiện spam, dự đoán tiểu đường, phát hiện ung thư, v.v.
\end{dinhnghia}

\subsection{Các loại hồi quy logistic}

\begin{ynghia}[title={Phân loại theo số lớp}]
Thông thường ``hồi quy logistic'' nghĩa là hồi quy logistic \textbf{nhị phân} với biến mục tiêu nhị phân; tuy nhiên biến mục tiêu có thể có thêm các dạng khác.
Theo số lớp của biến mục tiêu, hồi quy logistic chia thành các loại sau.
\end{ynghia}

\begin{dinhnghia}[title={Nhị phân (Binary hoặc Binomial)}]
Trong phân loại này, biến phụ thuộc chỉ có hai giá trị: 1 và 0.
Ví dụ: thành công/thất bại, có/không, thắng/thua.
\end{dinhnghia}

\begin{dinhnghia}[title={Đa thức (Multinomial)}]
Biến phụ thuộc có \textbf{3 lớp trở lên}, không có thứ tự hoặc không có ý nghĩa định lượng.
Ví dụ: ``Loại A'', ``Loại B'', ``Loại C''.
\end{dinhnghia}

\begin{dinhnghia}[title={Thứ bậc (Ordinal)}]
Biến phụ thuộc có \textbf{3 lớp trở lên có thứ tự} hoặc có ý nghĩa định lượng.
Ví dụ: ``kém'', ``tốt'', ``rất tốt'', ``xuất sắc''; mỗi hạng có thể gán điểm 0, 1, 2, 3.
\end{dinhnghia}

\subsection{Giả định hồi quy logistic}

\begin{luuy}[title={Giả định trước khi triển khai}]
Trước khi triển khai hồi quy logistic, cần lưu ý các giả định sau.
\begin{itemize}
  \item Với hồi quy logistic nhị phân, biến mục tiêu phải \textbf{luôn nhị phân}; kết quả mong muốn được biểu diễn bởi mức nhân tố 1.
  \item Mô hình \textbf{không} được có đa cộng tuyến (multi-collinearity): các biến độc lập phải độc lập với nhau.
  \item Chỉ đưa các biến \textbf{có ý nghĩa} vào mô hình.
  \item Nên chọn \textbf{cỡ mẫu lớn} cho hồi quy logistic.
\end{itemize}
\end{luuy}

\subsection{Mô hình hồi quy logistic nhị phân}

\begin{dinhnghia}[title={Dạng đơn giản nhất}]
Dạng đơn giản nhất là hồi quy logistic nhị phân (binomial): biến mục tiêu chỉ có hai giá trị 1 hoặc 0.
Cho phép mô hình hóa quan hệ giữa nhiều biến dự báo và biến mục tiêu nhị phân.
Trong hồi quy logistic, hàm tuyến tính được dùng làm đầu vào cho hàm khác:
\[
h_{\theta}(x) = g(\theta^{T} x)
\]
ở đây $g$ là hàm \textbf{logistic} hoặc \textbf{sigmoid}.
\end{dinhnghia}

\begin{dinhnghia}[title={Hàm sigmoid}]
\[
g(z) = \frac{1}{1 + e^{-z}}
\]
Đường cong sigmoid có giá trị trục $y$ nằm trong khoảng $(0, 1)$ và cắt trục tại $0{,}5$.
(Không có hình minh họa trong bộ tài liệu này; có thể vẽ $g(z)$ với $z \in [-6, 6]$ để thấy dạng chữ S.)
\end{dinhnghia}

\begin{phuongphap}[title={Phân lớp dương/âm}]
Các lớp có thể chia thành dương hoặc âm.
Đầu ra của hàm giả thuyết được hiểu là xác suất lớp dương nếu nằm trong $(0, 1)$.
Trong triển khai bài này: đầu ra $\geq 0{,}5$ $\Rightarrow$ dương; ngược lại $\Rightarrow$ âm.
\end{phuongphap}

\begin{dinhnghia}[title={Hàm mất mát $J(\theta)$}]
Cần định nghĩa hàm loss đo mức độ tốt của thuật toán với trọng số $\theta$:
\[
J(\theta) = \frac{1}{m}\left(-y^{T}\log(h) - (1 - y)^{T}\log(1 - h)\right)
\]
Mục tiêu chính: \textbf{tối thiểu hóa} hàm loss --- điều chỉnh trọng số (tăng/giảm).
Nhờ đạo hàm của loss theo từng trọng số, biết tham số nào nên lớn hoặc nhỏ hơn.
\end{dinhnghia}

\begin{dinhnghia}[title={Gradient descent}]
Phương trình gradient descent cho biết loss thay đổi thế nào khi sửa tham số:
\[
\frac{\partial J(\theta)}{\partial \theta_{j}} = \frac{1}{m} X^{T}(h - y)
\]
\end{dinhnghia}

\subsection{Triển khai hồi quy logistic nhị phân bằng Python}

\begin{dongco}[title={Dataset Iris (hai feature đầu)}]
Dùng bộ hoa iris: 3 lớp, mỗi lớp 50 mẫu; chỉ lấy \textbf{hai cột feature đầu}.
Mỗi lớp là một loài iris; chuyển bài toán thành phân loại nhị phân (lớp 0 vs các lớp khác).
\end{dongco}

\begin{vidu}[title={Import và nạp dữ liệu}]
\begin{lstlisting}
import numpy as np
import matplotlib.pyplot as plt
import seaborn as sns
from sklearn import datasets

iris = datasets.load_iris()
X = iris.data[:, :2]
y = (iris.target != 0) * 1
\end{lstlisting}
\end{vidu}

\begin{vidu}[title={Vẽ dữ liệu huấn luyện}]
\begin{lstlisting}
plt.figure(figsize=(6, 6))
plt.scatter(X[y == 0][:, 0], X[y == 0][:, 1],
            color='g', label='0')
plt.scatter(X[y == 1][:, 0], X[y == 1][:, 1],
            color='y', label='1')
plt.legend()
plt.show()
\end{lstlisting}
\end{vidu}

\begin{vidu}[title={Lớp LogisticRegression thủ công}]
Định nghĩa sigmoid, loss và gradient descent:

\begin{lstlisting}
class LogisticRegression:
    def __init__(self, lr=0.01, num_iter=100000,
                 fit_intercept=True, verbose=False):
        self.lr = lr
        self.num_iter = num_iter
        self.fit_intercept = fit_intercept
        self.verbose = verbose

    def __add_intercept(self, X):
        intercept = np.ones((X.shape[0], 1))
        return np.concatenate((intercept, X), axis=1)

    def __sigmoid(self, z):
        return 1 / (1 + np.exp(-z))

    def __loss(self, h, y):
        return (-y * np.log(h) - (1 - y) * np.log(1 - h)).mean()

    def fit(self, X, y):
        if self.fit_intercept:
            X = self.__add_intercept(X)
        self.theta = np.zeros(X.shape[1])
        for i in range(self.num_iter):
            z = np.dot(X, self.theta)
            h = self.__sigmoid(z)
            gradient = np.dot(X.T, (h - y)) / y.size
            self.theta -= self.lr * gradient
            z = np.dot(X, self.theta)
            h = self.__sigmoid(z)
            loss = self.__loss(h, y)
            if self.verbose and i % 10000 == 0:
                print(f'loss: {loss} \t')

    def predict_prob(self, X):
        if self.fit_intercept:
            X = self.__add_intercept(X)
        return self.__sigmoid(np.dot(X, self.theta))

    def predict(self, X):
        return self.predict_prob(X).round()
\end{lstlisting}
\end{vidu}

\begin{vidu}[title={Huấn luyện, dự đoán và đánh giá}]
\begin{lstlisting}
model = LogisticRegression(lr=0.1, num_iter=300000)
preds = model.predict(X)
(preds == y).mean()
\end{lstlisting}

Vẽ biên quyết định (đường contour xác suất 0.5):

\begin{lstlisting}
plt.figure(figsize=(10, 6))
plt.scatter(X[y == 0][:, 0], X[y == 0][:, 1],
            color='g', label='0')
plt.scatter(X[y == 1][:, 0], X[y == 1][:, 1],
            color='y', label='1')
plt.legend()
x1_min, x1_max = X[:, 0].min(), X[:, 0].max()
x2_min, x2_max = X[:, 1].min(), X[:, 1].max()
xx1, xx2 = np.meshgrid(np.linspace(x1_min, x1_max),
                       np.linspace(x2_min, x2_max))
grid = np.c_[xx1.ravel(), xx2.ravel()]
probs = model.predict_prob(grid).reshape(xx1.shape)
plt.contour(xx1, xx2, probs, [0.5], linewidths=1, colors='red')
plt.show()
\end{lstlisting}
\end{vidu}

\begin{ynghia}[title={Đánh giá mô hình nhị phân}]
Tỷ lệ \texttt{(preds == y).mean()} là độ chính xác trên toàn bộ tập đã dùng để fit (bài gốc).
Trong thực hành nên tách train/test để đánh giá không chệch.
\end{ynghia}

\subsection{Mô hình hồi quy logistic đa thức}

\begin{dinhnghia}[title={Multinomial logistic regression}]
Dạng hữu ích khác: hồi quy logistic \textbf{đa thức} --- biến mục tiêu có 3 lớp trở lên \textbf{không có thứ tự}, không có ý nghĩa định lượng.
\end{dinhnghia}

\subsection{Triển khai hồi quy logistic đa thức bằng Python}

\begin{dongco}[title={Dataset digits}]
Dùng tập \texttt{digits} từ scikit-learn (chữ số viết tay 0--9).
\end{dongco}

\begin{vidu}[title={Import thư viện}]
\begin{lstlisting}
from sklearn import datasets
from sklearn import linear_model
from sklearn import metrics
from sklearn.model_selection import train_test_split
\end{lstlisting}
\end{vidu}

\begin{vidu}[title={Nạp digits và chia tập}]
\begin{lstlisting}
digits = datasets.load_digits()
X = digits.data
y = digits.target
X_train, X_test, y_train, y_test = train_test_split(
    X, y, test_size=0.4, random_state=1)
\end{lstlisting}
\end{vidu}

\begin{vidu}[title={Huấn luyện và đánh giá}]
\begin{lstlisting}
digreg = linear_model.LogisticRegression()
digreg.fit(X_train, y_train)
y_pred = digreg.predict(X_test)
print("Accuracy of Logistic Regression model is:",
      metrics.accuracy_score(y_test, y_pred) * 100)
\end{lstlisting}

\textbf{Output:}
\begin{verbatim}
Accuracy of Logistic Regression model is: 95.6884561891516
\end{verbatim}
\end{vidu}

\begin{ynghia}[title={Kết quả độ chính xác}]
Độ chính xác mô hình khoảng \textbf{96\%} (cụ thể $\approx 95{,}69\%$ trên tập test với \texttt{test\_size=0.4}, \texttt{random\_state=1}).
\end{ynghia}
